% Section 1, part 2: the statements. Continued in part 3.

%-------------------------------------------------------------------------------
\subsection{Results}
\label{sec:int-results}

For a leg-free stable graph $\Gamma$ of genus $g$, a matrix series $X$ and a cohomological field theory $\Omega$ we write $\mathrm{Cont}_\Gamma[E_X,q_X;\Omega]$ for the contraction of the vertex tensors of $\Omega$ along the edges $E_X$ of $X$, with the leaf point $q_X$ of $X$, in the convention of \cite[Defs.~2.2, 2.13]{PPZ15}; it includes the factor $|\Aut\Gamma|^{-1}$ (\cref{def:set-graphsum,def:flh-outer}). The \emph{vertex graph} is the graph with one vertex, of genus $g$, and no edge; its term contains the leaf decorations.

\begin{thm}[Theorem~A]
\label{thm:int-main}
Let $g\geq2$ and let $(\alpha,\beta)\in\bbQ^2$ define a conifold frame (\cref{def:set-frames}), with vanishing period $\varpi_1=u+O(u^2)$ and $\tc=\varpi_1/\varpi_0$. Let $M:=\Mo(\tau(u)\zh)$, exactly in $u$, and $R':=\Qr M^{-1}$, where $\Qr$ is the matrix of the lowered rows of \cref{def:set-rows} expanded in the frame (\cref{def:set-confexp}). For a leg-free stable graph $\Gamma$ of genus $g$ let
\[
T_\Gamma:=\mathrm{Cont}_\Gamma\bigl[E_{R'},q_{R'};M.\TFT\bigr]
\]
be its outer term (\cref{def:set-outer}), where $\TFT$ is the topological field theory of the formal quintic (\cref{def:set-frob}). Then:
\begin{enumerate}
\item $T_\Gamma\in\bbQ\fps{u}$ for every $\Gamma$ with at least one edge;
\item for the vertex graph, $T_\Gamma\in c_g\tc^{2-2g}+\bbQ\fps{u}$, where
\[
c_g=\frac{(-1)^{g-1}B_{2g}}{2g(2g-2)\,5^{g-1}}=\frac{|B_{2g}|}{2g(2g-2)\,5^{g-1}} .
\]
\end{enumerate}
\end{thm}

\Cref{thm:int-main} is proved as \cref{thm:mai-main}; for one $g$ it is Property N($g$) of \cref{def:set-N}. The values are $c_2=1/1200$, $c_3=1/25200$, $c_4=1/180000$, $c_5=1/660000$ (\cref{rmk:set-cg}). The terms of the hatted outer sum, formed with the hatted data of \cref{def:set-hat} (\cref{def:set-outer}), are $G_\Gamma=u^{2g-2}T_\Gamma$ (\cref{lem:flh-hat}), so (i) says that $G_\Gamma\in u^{2g-2}\bbQ\fps{u}$ for every graph with an edge. The statement is about each graph separately. With (T) it implies the statement about their sum, Corollary~B below; it is not implied by a statement about the singular part of the sum: the computations recorded in \cref{sec:evi-univ} show that if the canonical gauge of the proof is replaced by a gauge in a non-canonical normalisation, six edge graphs of genus $4$ become singular while the singular part of the total is unchanged (\cref{rmk:mai-graphlocal}).

\begin{cor}[Corollary~B: the conifold gap of the formal quintic]
\label{cor:int-gap}
Let $g\geq2$ and $(\alpha,\beta)\in\bbQ^2$. The conifold expansion $F_g^{\mathrm{cf}}\in\bbQ\lser{s}$ of the genus-$g$ potential of the formal quintic in the frame $(\alpha,\beta)$ (\cref{def:mai-Fcf}) satisfies
\[
F_g^{\mathrm{cf}}\in c_g\tc^{2-2g}+\bbQ\fps{u} .
\]
Since $\bbQ\fps{u}=\bbQ\fps{\tc}$, the series $F_g^{\mathrm{cf}}$ has no fractional power of $u$, no term $\tc^k$ with $3-2g\leq k\leq-1$, and leading term $c_g\tc^{2-2g}$.
\end{cor}

Here $F_g^{\mathrm{cf}}$ is the image under $\Cexp_{\alpha,\beta}$ of the genus-$g$ graph sum written in the differential ring $\Nring$, whose evaluation is $F_g$ by Teleman's reconstruction and (QRR) (\cref{def:mai-Fcf}). \Cref{cor:int-gap} is proved as \cref{cor:mai-gap}: by the group property of the unit-preserving $R$-matrix action \cite[Def.~2.13, Prop.~2.14]{PPZ15}, $F_g^{\mathrm{cf}}$ is the sum of the outer terms $T_\Gamma$, and \cref{thm:int-main} applies to each of them.

The proof of \cref{thm:int-main} compares two normalisations of the $R$-matrix. The cyclic system of the formal quintic and that of the model at $v(u)$ are related by a gauge $P(u,z)$ that is regular at $u=0$ and unique up to a scalar factor $\cn(z)$ independent of $u$ (Theorem~U, \cref{thm:gau-U}); the \emph{canonical gauge} $\Pc$ is fixed at the conifold point by the gauge identity $\Pc\YC(v(u))=\YCQ$ between the vanishing solutions of the two systems (\cref{thm:gau-G,def:gau-can}). The actual rows $\Qr$ are normalised at $q=0$ by (QRR). Both $\Qr$ and $\Pc\Qr^{\mathrm m}(v(u))$, $\Qr^{\mathrm m}$ the model rows, are formal WKB solution matrices of the same system, with the same exponentials; they differ by a constant diagonal normalisation series,
\[
\Qr\,\Dn=\Pc\,\Qr^{\mathrm m}(v(u)),\qquad \Dn=\diag(\Dn_0,\dots,\Dn_4),\qquad \Dn_i(z)=\Dn_0(\zeta^{-i}z),\qquad \Dn_0\in\bbQ\fps{z}
\]
(\cref{thm:gau-R0,cor:gau-colours}; that $\Dn_0$ has rational coefficients is part of \cref{thm:con-glue}).

\begin{thm}[Theorem~C: the normalisation]
\label{thm:int-D}
$\Dn_0=1$ in $\bbQ\fps{z}$. Hence $\Dn=1$, and the window $\Fg{g}$, that is $\Dn_0\equiv1\bmod z^{3g-2}$ (\cref{def:gau-Fg}), holds for every $g\geq2$. Equivalently, in every conifold frame $\Qr=\Pc\,\Qr^{\mathrm m}(v(u))$: the actual gauge is the canonical gauge, and the $R$-matrix of the formal quintic, normalised at $q=0$ by (QRR), is in the canonical normalisation at the conifold point.
\end{thm}

\Cref{thm:int-D} is \cref{thm:gam-D} together with \cref{cor:mai-actual}; its consequences for the natural form are \cref{cor:mai-natfact}.

The identity $\Dn=1$ is reduced to one value at $q=0$ of a saddle series, the content of the next theorem.

\begin{thm}[Theorem~Gamma]
\label{thm:int-gamma}
Let $c(z;L)=\sum_{n\geq0}c_n(L)z^n$, $c_0(L)=1$, be the full-Stirling saddle series of \cref{def:gam-c}: the normalised stationary-phase series at $\sigma=L$ of
\[
\Gamma\Bigl(1+\frac{5\sigma}{z}\Bigr)\prod_{j=0}^{4}\Gamma\Bigl(\frac{\zeta^j-\sigma}{z}\Bigr)(-q)^{\sigma/z}\,\dd\sigma,\qquad q=q(L)=\frac{1-L^{-5}}{5^5},
\]
with every Gamma factor replaced by its formal Stirling form (\cref{def:fsp-stirling}). Then every $c_n$ lies in $\AL=\bbQ[L,L^{-1},(L^5-1)^{-1}]$, and:
\begin{enumerate}
\item[(a)] no $c_n$ has a pole at $L=1$, that is, at $q=0$, and $c(z;1)=h(z)$;
\item[(b)] no $c_n$ has a pole at $L=0$, that is, at $q=\infty$, and
\[
c(z;0)=\Ginf(z):=\exp\Bigl(\sum_{l\geq1}\frac{(-1)^{l+1}B_{5l+1}(\frac15)}{l(5l+1)}\,z^{5l}\Bigr),
\]
the series given by \cite[Lemma~7.7, (37)]{GJR18} at $m=5$, $k=1$, $\lambda=1$ in the notation there.
\end{enumerate}
\end{thm}

\Cref{thm:int-gamma} is formed by parts (i), (ii) and the first assertion of (iv) of \cref{thm:gam-main}, with \cref{thm:gam-ginf}. Part (a) is proved by a local Mellin inversion (\cref{thm:gam-lmi}), from the formal stationary-phase lemmas of \cref{sec:fsp} alone. The formula in (b) is the one of \cite[Lemma~7.7, (37)]{GJR18}, proved there for the Landau--Ginzburg series of \cite{GJR18}; the content of (b) here is that the series $c(z;L)$ takes this value at $L=0$, and we do not claim that the series of \cite{GJR18} coincides with $c(z;L)$ (\cref{thm:gam-ginf}). Part (b) identifies the free coefficients of the canonical colour-$0$ function (\cref{cor:gam-free}); it is not used in the proof of \cref{thm:int-main}. By the identity $c_n(L)=L^{4n}\mathrm{ph}^{\mathrm{SP}}_n(L^{-5})$ between the flat and the hatted saddle series (\cref{prop:gam-hatflat}), Theorem~P (\cref{thm:con-Pb}) and the glue (\cref{thm:con-glue}), part (a) is equivalent to $\Dn=1$ (\cref{cor:gam-equiv,cor:con-A2}); this is how \cref{thm:int-D} is proved.

\begin{rmk}
\label{rmk:int-constant}
The constant $c_g$ is the constant of \cite[Conjecture~1.5]{CGYZ26} for the quintic threefold, with the same sign and the same power of $5$, in the same normalisation of the vanishing period: the formal quintic has the Picard--Fuchs operator $\PFfour$ of the quintic, the period $\omega_1=\Delta+\frac7{10}\Delta^2+O(\Delta^3)$ of \cite[Example~1.3]{CGYZ26} is $\varpi_1$ (with $\Delta=u$), the pair of \cite[Example~1.3]{CGYZ26} is the frame $(\alpha,\beta)=(0,0)$, and the family of \cite[Remark~1.4]{CGYZ26} with rational coefficients is the family of \cref{def:set-frames}; the normalisation $\omega_1=\Delta+O(\Delta^2)$ is the one of \cite[Example~1.3; \S4.2, Proposition-Definition~4.5 (1)]{CGYZ26}. The conifold expansion of \cref{def:set-confexp} and that of \cite[Definition~1.2]{CGYZ26} agree on the differential ring generated by $C_0^{\pm1}$ and $C_1^{\pm1}$ over $\bbQ[L^5,L^{-5}]$, wherever the latter is defined. The two theories have different potentials, so the equality of the constants compares two different series, and nothing about the quintic threefold follows from it. The dictionary and the differences are \cref{rmk:mai-constant}. \draftnote{Numbers from arXiv:2609.13921v2, read in the HTML rendering; to be re-read in the PDF.}
\end{rmk}

% Section 1, part 3: inputs, what is not claimed, strategy, plan of the note.

%-------------------------------------------------------------------------------
\subsection{Inputs}
\label{sec:int-inputs}

The results are conditional on the following inputs. Statement numbers refer to the arXiv versions named in the bibliography; each input is stated again at the place where it enters.
\begin{description}
\item[(Q)] \emph{Quoted theorems.} Topological recursion, applied only to the model curve $x=\Wv(p)$, $y=p$ (\cref{sec:mod}): \cite[Def.~4.2, (4-9), with the kernel (4-3); Theorems~4.2 and~4.6; Theorem~5.1, (5-14), with \S5.3, (5-27)--(5-30)]{Eynard:2007kz}; the theorem \cite[Theorem~4.1]{DuninBarkowski:2012bw}, with its graph weights \cite[\S2.3]{DuninBarkowski:2012bw}, as recalled in \cite[\S2.6, (2-30)--(2-35)]{Dunin-Barkowski:2015caa}, with the action of \cite[\S2.3]{Dunin-Barkowski:2015caa} (its dilaton leaves taken with the sign of \cite[Def.~2.13]{PPZ15}), the compatibility condition \cite[Prop.~4.3 (1)]{Dunin-Barkowski:2015caa} and the sentence preceding \cite[Theorem~8.1]{Dunin-Barkowski:2015caa}, applied to $f=4\Wv$, that is, to \cite[(8-1)]{Dunin-Barkowski:2015caa} with $n=4$ and $a=(0,0,0,0,4v)$, whose critical points $p^5=v$ have the pairwise distinct critical values $5p_c^4$; and \cite[Theorem~7.1: (7.6), (7.7) with its factor $1/\hbar^2$, (7.9), (7.10), with the kernel (1.1)]{ABDKS25}. The $R$-matrix action: \cite[\S0.2, (1); Def.~0.3 (iii); (5); Remark~2.1 and (6); Def.~2.2, with the convention stated after it; Prop.~2.4; Def.~2.5, with the convention stated after it; Props.~2.8 and~2.9; Prop.~2.12, with the formula (8) in its proof; Def.~2.13; Prop.~2.14]{PPZ15}; Props.~2.4, 2.8 and~2.9 enter only through the proof of Prop.~2.14, and (8) only through \cref{lem:set-pushins}.
\item[(QRR)] \emph{The value of the $R$-matrix at the large-radius point}: $R(q=0)=\diag(h_{\bw_0},\dots,h_{\bw_4})$ in the normalised canonical basis at $q=0$, in Givental's convention (the columns $R\,\tilde e_i\,\re^{u_i/z}$ solve the quantum differential equation, and the legs of the graph sum carry $R^{-1}(\psi)$) (\cref{thm:set-qrr}). It is used for every colour, in the form $\Phi_i(z,1)=h_{\bw_i}(z)$ (\cref{eq:set-qrr}): in \cref{def:set-rows}, in \cref{lem:gau-colrows} and, for colour $0$, in the glue \cref{thm:con-glue}. Sources: \cite[(17), in the proof of Lemma~5.2]{GJR18}; \cite[(19)--(20), in the proof of Lemma~4.6, at $t=0$]{GJR17}; \cite[\S2.2, Theorem~1 and the Remark after it]{Giv01}; \cite[Theorem~1, Corollary~1]{MR2276766}. How (QRR) follows from the cohomological field theory is explained in \cref{rmk:set-qrrderiv}, which no proof uses.
\item[(M)] \emph{The genus-zero mirror theorem} \cite[Theorem~11.6]{MR1408320}\draftnote{numbers from arXiv v2; IMRN numbering not compared}, in the Calabi--Yau case $n=4$, $r=1$, $l_1=5$, $\lambda'=0$, with the set-up of \cite[\S9]{MR1408320} and \cite[Prop.~11.4]{MR1408320}: at $\lambda_i=\zeta^i$, $J(T(q)H,z)=I(q,z)/C_0(q)$ with $T=I_1/C_0$ (\cref{thm:set-mirror}).
\item[(G0)] \emph{The genus-zero axioms of the twisted theory}: the genus-zero twisted classes form a cohomological field theory with unit for $\eta$ \cite[\S4.1]{GJR18}, \cite[Def.~0.3]{PPZ15}, polynomial and $S_5$-symmetric in $\lambda$, with the string, divisor and splitting axioms, the genus-zero topological recursion relation and the quantum differential equation \cite[\S5.2, (14)]{GJR18} (\cref{thm:set-G0}).
\item[(T)] \emph{Teleman's reconstruction} \cite[Theorem~(2) in (1.15)]{MR2917177} with \cite[\S4.1, Theorem~4.1]{GJR18}: the cohomological field theory of the formal quintic is the graph sum of \cref{def:set-graphsum} with the $R$-matrix of \cref{thm:set-qrr} (\cref{thm:set-GT}). In \cref{thm:int-main,thm:int-D} it identifies that $R$-matrix, whose rows define $T_\Gamma$ and $\Dn$; in \cref{cor:int-gap} it identifies the total with $F_g$.
\end{description}
In addition, the formula \cite[(1.7)]{AC96}, with $\kappa_0=2g-2+n$ \cite[p.~5]{AC96}, is used in part (i) of the dilaton lemma \cref{lem:set-dilaton}. The definitions use \cite[\S0.2, (5)]{LP18}, \cite[\S0.2, (5), (10)]{LP20}, \cite{ZZ08} and \cite{MR1666787}. The inputs are read with the following definitions and conventions of their sources: the normalisation of the $J$-function \cite[Section~6, (10)]{MR2276766} in (M) (\cref{thm:set-mirror}); the free energies \cite[Def.~4.3, (4-10)]{Eynard:2007kz} (\cref{conv:mod-kernel,cor:mod-trF}) and the set-up \cite[(5-1), (5-4), (5-13)]{Eynard:2007kz} of \cite[Theorem~5.1]{Eynard:2007kz} (\cref{thm:mod-var}); and the prescription \cite[\S3.1, \S3.5]{DuninBarkowski:2012bw}, by which the integrals \cite[(2-31), (2-32)]{Dunin-Barkowski:2015caa} are read as formal Laplace series (\cref{lem:mod-laplace,prop:mod-Rtr}).

\emph{Where the inputs enter.} (M) and (G0) enter only in \cref{sec:set}, through the genus-zero data and the large-radius cyclic structure (\cref{thm:set-mirror,prop:set-cyclicLR,cor:set-cyclicPF5}). (QRR) enters at the three places named above; colour $0$ of the comparison of the actual rows with the canonical gauge (\cref{thm:gau-R0}) uses no (QRR). (T) enters in \cref{thm:set-GT} and, for \cref{cor:int-gap}, in \cref{cor:mai-gap}. The quoted results of \cite{PPZ15} enter in \cref{sec:set} (\cref{thm:set-G0}, as part of (G0); \cref{def:set-graphsum,lem:set-pushins,def:set-modcohft,lem:set-modhom,lem:set-dilaton,def:set-outer}), in \cref{sec:tra-dress}, in \cref{sec:mod} (\cref{prop:mod-Rtr}, through \cref{lem:set-pushins}, and \cref{thm:mod-tr}), in \cref{thm:flh-main} and in \cref{cor:mai-gap}; the formula of \cite{AC96} enters through \cref{lem:set-dilaton}. Topological recursion enters only in \cref{sec:mod}, where the Model Lemma is proved; nothing in \cref{sec:set,sec:gau,sec:tra,sec:red} or \cref{sec:flh,sec:con,sec:gam} uses it. In \cref{sec:gau} the cyclic form of $\PFfour$, Theorem~U, the gauge identity and symplecticity use no quoted input beyond the genus-zero data of \cref{sec:set}; the transport lemmas of \cref{sec:tra} use \cite[Defs.~2.2, 2.13]{PPZ15} and the dilaton lemma; Theorem~R adds no input; the comparison theorem of \cref{sec:flh} uses \cite[Def.~0.3, Remark~2.1 and (6), Defs.~2.2, 2.5, 2.13, Prop.~2.12]{PPZ15} and the dilaton lemma; Theorem~P(a) (\cref{thm:con-Pa}) and part (a) of \cref{thm:int-gamma} use no quoted input.

\emph{Conditionality.} Theorem~A holds modulo the quoted theorems (Q) of \cite{Eynard:2007kz}, \cite{DuninBarkowski:2012bw} (as recalled in \cite{Dunin-Barkowski:2015caa}), \cite{ABDKS25} and \cite{PPZ15} with the statement numbers above, the value (QRR) of the $R$-matrix at the large-radius point, the genus-$0$ mirror theorem (M) and the genus-$0$ axioms (G0) of the twisted theory, Teleman's reconstruction (T), and the formula \cite[(1.7)]{AC96}; every other step is proved in the note.

In particular Theorem~U (over every field of characteristic $0$, \cref{thm:gau-U}) and every lemma of the chain are proved in the note; Stirling's series enters only through its formal proof as the stationary phase of Euler's integral (\cref{lem:fsp-stirling}), and the implicit function theorem only in its formal version (\cref{lem:fsp-ift}). \Cref{cor:int-gap} uses, beyond \cref{thm:int-main}, only (T), \cite[\S0.2, (1); Def.~0.3; (5); Remark~2.1; Defs.~2.2 and~2.5, with the conventions stated after them; (8); Prop.~2.12; Def.~2.13; Prop.~2.14, with Props.~2.4, 2.8 and~2.9 through its proof]{PPZ15} and statements proved in the note, among them the comparison of the two forms of the translation (\cref{lem:set-pushins}). \Cref{thm:int-D}, like \cref{thm:int-main}, concerns the rows of the $R$-matrix that (T) identifies (\cref{thm:set-GT}); its proof uses this matrix only through \cref{def:set-rows}~(b), $R(0)=1$ and (QRR), the last through the glue \cref{thm:con-glue} and, for the colours $i\neq0$, through \cref{lem:gau-colrows} (by \cref{cor:gau-colours}), and it uses the genus-zero data of \cref{sec:set}, which rest on (M) and (G0), and the results of \cref{sec:gau,sec:con,sec:gam} and of \cref{sec:fsp} (\cref{thm:gam-D,cor:mai-actual}). \Cref{thm:int-gamma} uses no quoted input, except that its part (b) uses the identity \cite[24.4.1]{DLMF} of Bernoulli polynomials (\cref{thm:gam-ginf}). No finite computation enters a proof: the computations that accompanied the work are collected in \cref{sec:evi}, and no proof refers to them.

Apart from the inputs above, the citations in definitions and the identity \cite[24.4.1]{DLMF} used for part (b) of \cref{thm:int-gamma}, the citations of the note are further sources of the inputs, comparisons with the literature, conventions (\cref{sec:sgn}) or context, and they enter no proof. In particular the following are cited but are not used in the proofs of \cref{thm:int-main,cor:int-gap}: \cite[Corollary~1.2]{Hock25}, which gives a first proof of the Gap Theorem (\cref{rmk:mod-hock}); \cite[Theorem~4.4, (4-23)--(4-24); Corollary~4.1, (4-26)--(4-29); (A-25)]{Eynard:2007kz}, which \cref{lem:mod-xfun} does not need; the dilaton equation \cite[Theorem~4.7, (4-30)]{Eynard:2007kz} and the variation (5-15) of the free energies in \cite[Theorem~5.1]{Eynard:2007kz}, recalled in \cref{cor:mod-trF} and \cref{thm:mod-var}~(ii); Mumford's relation in the form of \cite[\S2.1]{MR1728879}, \cite[(26)]{LP20} and the statements of \cite[\S8]{MR2917177} cited in \cref{rmk:set-qrrderiv}, which enter only \cref{rmk:set-qrrsign} and \cref{rmk:set-qrrderiv}; \cite[Theorem~2]{MR2276766}, a modern formulation of (M), and \cite[Corollary~11.7]{MR1408320}, a misprinted form of it; \cite[Props.~6.1--6.2]{GR19} and \cite[\S2.4.3, (17); \S2.4.4, (19)--(20)]{LP20}, which state relations equivalent to \cref{prop:set-cyclicLR}; and \cite{CGYZ26}, in \cref{rmk:int-constant,rmk:mai-constant}.

%-------------------------------------------------------------------------------
\subsection{What is not claimed}
\label{sec:int-notclaimed}

\begin{enumerate}
\item Nothing is claimed about the quintic threefold or about its graph sums. The formal quintic is a different theory (already its genus-$2$ constant-map term differs from that of the quintic, \cref{sec:evi-calib}), and the agreement of the constants in \cref{rmk:int-constant} compares two different potentials. The conifold gap of the quintic threefold is \cite[Theorem~2]{CGYZ26}.
\item Nothing is claimed for frames other than those of \cref{def:set-frames} and the rescalings of their two periods treated in \cref{rmk:mai-constant}, for $g=1$, or away from the conifold point $u=0$.
\item (QRR) is quoted, not proved; \cref{rmk:set-qrrderiv} explains how it follows from the cohomological field theory and is not used.
\item The second key statement of \cite[Remark~1.9]{CGYZ26}, the conifold regularity of correlators of the formal quintic with special insertions, is not addressed: the Model Lemma (\cref{def:red-ML}) is a statement about the model $\Wv$ and is not that statement. Nor is \cref{cor:int-gap} claimed to be literally the first key statement of \cite[Remark~1.9]{CGYZ26}, which is named there but not stated.
\item No novelty is claimed for the formal stationary-phase lemmas of \cref{sec:fsp}, which are standard in spirit and are proved for completeness, for the formula of $\Ginf$ \cite[Lemma~7.7, (37)]{GJR18}, or for the change of chart \cite[\S10]{MR1901075}. For \cref{thm:int-gamma} the note cites \cite[Remark~7.2]{GJR18}, where the analogous identity for the Landau--Ginzburg integral of Guo, Janda and Ruan is stated without proof, and we do not claim that the Landau--Ginzburg series of \cite{GJR18} coincides with $c(z;L)$.
\item Topological recursion is used only on the model curve $x=\Wv(p)$, $y=p$. It is not applied to a mirror curve of the quintic or of the formal quintic, and \cite{FCCZ26} is not used.
\end{enumerate}

%-------------------------------------------------------------------------------
\subsection{Strategy}
\label{sec:int-strategy}

\emph{Reduction to the model (\cref{sec:gau,sec:tra,sec:red}).} In a conifold frame the rows of the formal quintic satisfy a cyclic system $z\,\de_uY_a=b_aY_{a+1}$, $a\in\bbZ/5$, with one simple pole at $u=0$: the four forward relations are the large-radius cyclic structure (\cref{prop:set-cyclicLR}), which rests on (M) and (G0), carried to the conifold point by the conifold expansion (\cref{prop:gau-quintic}), and the closing relation comes from the cyclic form of $\PFfour$, proved in every frame (\cref{thm:gau-cyc,cor:gau-cycsys}). The model, pulled back by $v(u)=u\tau(u)^{-5/4}$, gives a second such system with the same product of coefficients, and this equality of products is what fixes $\tau$ (\cref{prop:gau-equalprod}). Theorem~U (\cref{thm:gau-U}) gives a gauge between the two systems that is regular at $u=0$ and unique up to a scalar factor $\cn(z)$ independent of $u$; the gauge identity (\cref{thm:gau-G}) singles out the canonical gauge $\Pc$, which is symplectic (\cref{prop:gau-symp}). The vanishing solution $\YC$ of the model defines a tower of covectors $\Pi_i$ and a slice $\sK$ (\cref{def:tra-tower,def:tra-slice}); the edges of $\Pc$ lie in $\sKo\tensor\sKo$, its leaf point, rescaled by the dilaton factor, lies in $-z+\sK$, and every graph term is a finite sum (\cref{prop:tra-can}). Theorem~R (\cref{thm:red-main,cor:red-can}) reduces the regularity of every term with an edge of the flat graph sum of $\Pc.\Omega^{\mathrm{mod}}_{v(u)}$, and the vertex law, to the Model Lemma on $\sK$ (\cref{def:red-ML}), a statement about the model alone.

\emph{The model (\cref{sec:mod}).} The model cohomological field theory is identified with topological recursion on the curve $x=\Wv(p)$, $y=p$ (\cref{thm:mod-tr}), the $R$-matrix of the curve being $\Mo$ (\cref{prop:mod-Rtr}). Insertions from $\sK$ become legs on the curve (\cref{prop:mod-LD,thm:mod-lattice}), their pairings are polynomial in $v$ (\cref{thm:mod-Lreg,thm:mod-V}), and the free energies are $c_hv^{2-2h}$ (\cref{thm:mod-gap}); together these give the Model Lemma in every genus (\cref{thm:mod-MLK}).

\emph{The normalisation (\cref{sec:flh,sec:con,sec:gam}).} The comparison theorem (\cref{thm:flh-main}) identifies the outer term $T_\Gamma$ of \cref{thm:int-main} with the flat term of the reduction, graph by graph, provided $\Dn_0\equiv1\bmod z^{3g-2}$. The series $\Dn$ is a connection constant between the conifold point, where $\Pc$ is normalised, and $q=0$, where (QRR) normalises the actual rows. Theorem~P (\cref{thm:con-Pa,thm:con-Pb}) identifies the canonical colour-$0$ WKB function with a full-Stirling stationary-phase series $\PhiSP$, and the glue reads (QRR) as $h(z)\Dn_0(z)=\Phican(z,1)$ (\cref{thm:con-glue}); so $\Dn=1$ is the identity $\PhiSP(z,1)=h(z)$. Through $c_n(L)=L^{4n}\mathrm{ph}^{\mathrm{SP}}_n(L^{-5})$ (\cref{prop:gam-hatflat}) this is part (a) of \cref{thm:int-gamma}, proved by local Mellin inversion (\cref{thm:gam-lmi}). Hence $\Dn=1$ (\cref{thm:gam-D}), and \cref{thm:int-main} follows by assembling the three parts (\cref{thm:mai-main}).

%-------------------------------------------------------------------------------
\subsection{Plan of the note}
\label{sec:int-plan}

\Cref{sec:set} defines the formal quintic and its genus-zero data, states the inputs (M), (G0), (QRR) and (T), proves the large-radius cyclic structure, and defines the conifold frames and the conifold expansion, the model, $\Mo$, the model cohomological field theory and the outer graph sum. \Cref{sec:gau} proves the cyclic form of $\PFfour$, Theorem~U, the gauge identity and symplecticity, and compares the actual rows with the canonical gauge. \Cref{sec:tra} introduces the vanishing tower and the slice and proves the transport lemmas, and \cref{sec:red} proves Theorem~R. \Cref{sec:mod} proves the Model Lemma by topological recursion on the model curve. \Cref{sec:flh} proves the comparison of the flat and the outer graph sums. \Cref{sec:con} defines the saddle series $c(z;L)$ and proves Theorem~P and the glue, and \cref{sec:gam} proves \cref{thm:int-gamma} and \cref{thm:int-D}. \Cref{sec:mai} assembles \cref{thm:int-main}, proves \cref{cor:int-gap} and records the natural factorisation and universality. \Cref{sec:fsp} proves the lemmas of formal stationary phase, \cref{sec:sgn} collects the sign conventions and normalisations, and \cref{sec:evi} records the finite computations; nothing in the body refers to \cref{sec:evi} for a proof.

